\section{Optimization passes}

Canonicalization and constant folding run through the built in greedy driver,
using the folders and canonicalization patterns attached to the ops. Constant
subgraphs collapse to a single constant, and structural patterns such as relu
idempotence and reshape of reshape fire.

The BatchNorm folding pass is dedicated and does real tensor arithmetic at compile
time. At inference the scale, offset, mean, and variance of a batch normalization
are constants, so a batch norm following a convolution with no activation between
them is exactly another convolution with rescaled weights and bias:
\[
  \mathrm{coef}[o] = \frac{\mathrm{scale}[o]}{\sqrt{\mathrm{var}[o] + \epsilon}},
  \quad W'[o] = W[o]\,\mathrm{coef}[o],
  \quad b'[o] = \mathrm{coef}[o]\,(b[o] - \mathrm{mean}[o]) + \mathrm{offset}[o].
\]
The pass matches this pattern, computes the new constant tensors, and emits a
single convolution.

Operator fusion folds a trailing relu into a producing convolution or matmul by
setting the fused activation attribute. Because those ops already carry the bias
as an operand, the result is a single fused convolution or matmul that computes
the linear operation, the bias, and the activation together. The intermediate
therefore stays in the scratchpad rather than being written to DRAM and read back.

Dead code elimination needs no custom pass. Because every op is \texttt{Pure}, the
canonicalizer removes unused ops and \texttt{symbol-dce} removes unused private
functions. The three optimization levels compose these: \texttt{-O0} imports and
verifies, \texttt{-O1} adds canonicalization and folding, and \texttt{-O2} adds
fusion and dead code elimination.
